\section{Related Work and Theoretical Background}

\subsection{Agent-Based Computational Economics and Market Herding}
Agent-based computational economics provides a structural framework for modeling macroscopic market dynamics from heterogeneous agent interactions \cite{bouchaud2013more, farmer2009economy}. Early models demonstrated that stochastic recruitment dynamics generate clustering and bimodal sentiment states \cite{kirman1993ants}. Multi-agent market simulations subsequently showed that interaction between fundamentalist and chartist traders reproduces realistic power-law tails and volatility clustering \cite{lux1999scaling, alfarano2005estimation}. Further physical analogies, including spin-glass models with market-wide feedback, revealed that competing time horizons induce expectation bubbles and intermittency \cite{bornholdt2001expectation}.

\subsection{Informational Memory Queues in Swarm Physics}
In non-equilibrium statistical mechanics, autonomous agent populations with short-term cognitive memory exhibit spontaneous symmetry breaking \cite{li2026informational}. When individual agents record observations in discrete First-In-First-Out (FIFO) queues and compute decision postures via majority or linear score functions, the collective crowd state undergoes a phase transition at critical memory depth $M_c$. Below $M_c$, the crowd remains ergodic with equal probability of positive and negative consensus. Above $M_c$, non-linear peer observation feedback drives the crowd into stable polarized states. This paper adapts this physical formalism to real-world asset returns, demonstrating that financial traders operate as memory-constrained agents within a noisy observation channel.

\subsection{Backtest Overfitting and Statistical Significance in Quantitative Finance}
A persistent challenge in empirical quantitative finance is the vulnerability of backtested trading strategies to false discoveries \cite{bailey2017probability}. When multiple strategy configurations are evaluated on historical asset prices, standard in-sample optimization selects configurations that overfit past noise, leading to negative performance decay out-of-sample. Existing statistical approaches mitigate overfitting through multiple-testing adjustments (e.g., False Discovery Rate control and deflated Sharpe ratios). The PIMS framework introduces a physical complement to these statistical techniques by determining whether a strategy parameter configuration resides in an ergodic sub-critical regime ($k > 0$) or in a hysteresis-trapped super-critical regime ($k < 0$).

\subsection{Critical Transitions and Early-Warning Indicators}
Complex ecological and physical systems approaching tipping points exhibit characteristic dynamic slowing down, where recovery rates from external perturbations approach zero \cite{scheffer2009early}. In financial contexts, critical slowing down manifests during market panic events when crowd sentiment recovery times diverge exponentially. We formalize this phenomenon using relaxation half-life metrics ($T_{1/2}$) and prove that structural macro trend anchors provide an effective damping mechanism against permanent consensus lock-in.
